% ═══════════════════════════════════════════════════════════════
% LH-04b: El horario (LEHRERHINWEISE)
% Spanisch Klasse 7 - Sequenz 4: Mi colegio
% ═══════════════════════════════════════════════════════════════
% Abgeleitet von: LE-04b.tex (Score: 9.3/10)
% ═══════════════════════════════════════════════════════════════

\documentclass[11pt, a4paper]{article}

% ═══════════════════════════════════════════════════════════════
% SW-MODUS TOGGLE
% ═══════════════════════════════════════════════════════════════
\newif\ifswmodus
\swmodusfalse

% ═══════════════════════════════════════════════════════════════
% PAKETE
% ═══════════════════════════════════════════════════════════════

\usepackage[utf8]{inputenc}
\usepackage[T1]{fontenc}
\usepackage[ngerman]{babel}
\usepackage{helvet}
\renewcommand{\familydefault}{\sfdefault}

\usepackage[margin=2cm]{geometry}
\usepackage{enumitem}
\usepackage{tabularx}
\usepackage{booktabs}
\usepackage{longtable}

\usepackage{xcolor}
\usepackage{amssymb}
\usepackage{tcolorbox}
\tcbuselibrary{skins}

\usepackage{fancyhdr}
\usepackage{lastpage}
\usepackage{needspace}

% ═══════════════════════════════════════════════════════════════
% FARBEN
% ═══════════════════════════════════════════════════════════════

\definecolor{spanisch-color}{HTML}{c41e3a}
\definecolor{grau-color}{HTML}{888888}
\definecolor{hellgrau-color}{HTML}{f5f5f5}
\definecolor{basis-color}{HTML}{2a7a4b}
\definecolor{standard-color}{HTML}{2c5aa0}
\definecolor{erweiterung-color}{HTML}{7b1fa2}
\definecolor{loesung-color}{HTML}{c62828}

\definecolor{sw-dark}{HTML}{000000}
\definecolor{sw-gray}{HTML}{666666}
\definecolor{sw-white}{HTML}{ffffff}

\ifswmodus
    \colorlet{spanisch}{sw-dark}
    \colorlet{grau}{sw-gray}
    \colorlet{hellgrau}{sw-white}
    \colorlet{basis}{sw-dark}
    \colorlet{standard}{sw-dark}
    \colorlet{erweiterung}{sw-dark}
    \colorlet{loesung}{sw-dark}
\else
    \colorlet{spanisch}{spanisch-color}
    \colorlet{grau}{grau-color}
    \colorlet{hellgrau}{hellgrau-color}
    \colorlet{basis}{basis-color}
    \colorlet{standard}{standard-color}
    \colorlet{erweiterung}{erweiterung-color}
    \colorlet{loesung}{loesung-color}
\fi

\newcommand{\fachfarbe}{spanisch}

% ═══════════════════════════════════════════════════════════════
% VARIABLEN
% ═══════════════════════════════════════════════════════════════

\newcommand{\thema}{El horario -- Stundenplan und Uhrzeiten}
\newcommand{\sequenz}{04}
\newcommand{\block}{b}
\newcommand{\fach}{Spanisch}
\newcommand{\klasse}{7}
\newcommand{\dauer}{90 Min. (Doppelstunde)}
\newcommand{\lerngruppengroesse}{24}
\newcommand{\datum}{\today}

% ═══════════════════════════════════════════════════════════════
% KOPF- UND FUSSZEILE
% ═══════════════════════════════════════════════════════════════

\pagestyle{fancy}
\fancyhf{}
\renewcommand{\headrulewidth}{0.4pt}
\renewcommand{\footrulewidth}{0.4pt}

\lhead{\fach{} -- Klasse \klasse}
\chead{\textbf{LH-\sequenz\block{} (Lehrerhinweise)}}
\rhead{\datum}

\lfoot{}
\cfoot{}
\rfoot{Seite \thepage\ von \pageref{LastPage}}

% ═══════════════════════════════════════════════════════════════
% BEFEHLE
% ═══════════════════════════════════════════════════════════════

\newcommand{\B}{\textcolor{basis}{\textbf{[B]}}}
\newcommand{\Std}{\textcolor{standard}{\textbf{[S]}}}
\newcommand{\E}{\textcolor{erweiterung}{\textbf{[E]}}}
\newcommand{\loesung}[1]{\textcolor{loesung}{\textbf{#1}}}
\newcommand{\es}[1]{\textcolor{\fachfarbe}{\textbf{#1}}}

\newtcolorbox{kurzplan}{
    colback=hellgrau,
    colframe=\fachfarbe,
    colbacktitle=white,
    coltitle=\fachfarbe,
    boxrule=1pt,
    arc=0pt,
    fonttitle=\bfseries,
    attach boxed title to top left={yshift=-2mm, xshift=5mm},
    boxed title style={boxrule=0pt, colback=white},
    title=Kurzplan
}

\newtcolorbox{differenzierung}{
    colback=white,
    colframe=grau,
    colbacktitle=white,
    coltitle=grau,
    boxrule=0.5pt,
    arc=0pt,
    fonttitle=\bfseries,
    attach boxed title to top left={yshift=-2mm, xshift=5mm},
    boxed title style={boxrule=0pt, colback=white},
    title=Differenzierung
}

\newtcolorbox{fehlerbox}{
    colback=white,
    colframe=loesung,
    colbacktitle=white,
    coltitle=loesung,
    boxrule=0.5pt,
    arc=0pt,
    fonttitle=\bfseries,
    attach boxed title to top left={yshift=-2mm, xshift=5mm},
    boxed title style={boxrule=0pt, colback=white},
    title=Häufige Fehler
}

\newtcolorbox{materialbox}{
    colback=white,
    colframe=\fachfarbe,
    colbacktitle=white,
    coltitle=\fachfarbe,
    boxrule=0.5pt,
    arc=0pt,
    fonttitle=\bfseries,
    attach boxed title to top left={yshift=-2mm, xshift=5mm},
    boxed title style={boxrule=0pt, colback=white},
    title=Material-Checkliste
}

% ═══════════════════════════════════════════════════════════════
% DOKUMENT
% ═══════════════════════════════════════════════════════════════

\begin{document}

% ═══════════════════════════════════════════════════════════════
% KOPFBEREICH
% ═══════════════════════════════════════════════════════════════

\begin{center}
    {\Large\bfseries\textcolor{\fachfarbe}{Lehrerhinweise: \thema}}\\[0.3em]
    {\normalsize LH-\sequenz\block{} | \dauer}
\end{center}

\vspace{10pt}

% ═══════════════════════════════════════════════════════════════
% 1. KURZPLAN
% ═══════════════════════════════════════════════════════════════

\section*{1. Stundenverlauf (Kurzplan)}

\textbf{1. Stunde (45 Min.) -- Schwerpunkt: Wochentage + Uhrzeiten}

\begin{tabularx}{\textwidth}{|c|l|X|l|}
\hline
\textbf{Zeit} & \textbf{Phase} & \textbf{Inhalt} & \textbf{Material} \\
\hline
0--5 & Einstieg & Kalender: ``¿Qué día es hoy?'' & Kalender \\
\hline
5--15 & Input 1 & Wochentage einführen, vorsprechen & Tafel \\
\hline
15--20 & Übung 1 & Wochentage-Memory / Zuordnung & Karten \\
\hline
20--30 & Input 2 & Uhrzeiten mit Uhr: ``Son las...'' & Uhr \\
\hline
30--40 & Übung 2 & AB-04b Aufg. 1+2 & AB \\
\hline
40--45 & Sicherung & Lösungen besprechen & Tafel \\
\hline
\end{tabularx}

\vspace{1em}

\textbf{2. Stunde (45 Min.) -- Schwerpunkt: Stundenplan + Dialog}

\begin{tabularx}{\textwidth}{|c|l|X|l|}
\hline
\textbf{Zeit} & \textbf{Phase} & \textbf{Inhalt} & \textbf{Material} \\
\hline
0--5 & Aktivierung & Schnelles Uhrzeiten-Quiz & Uhr \\
\hline
5--15 & Input & Fragen: ¿A qué hora...? ¿Cuándo...? & Tafel \\
\hline
15--25 & Übung 1 & AB-04b Aufg. 3: Stundenplan & AB \\
\hline
25--35 & Übung 2 & PA: Dialog ``¿Cuándo tienes...?'' & -- \\
\hline
35--40 & Sicherung & AB-04b Aufg. 4: Dialog schreiben & AB \\
\hline
40--45 & Abschluss & Zusammenfassung, HA & -- \\
\hline
\end{tabularx}

% ═══════════════════════════════════════════════════════════════
% 2. DIFFERENZIERUNG
% ═══════════════════════════════════════════════════════════════

\section*{2. Differenzierung}

\begin{differenzierung}
\textbf{Legende:} \B{} Basis (BR) | \Std{} Standard (MR) | \E{} Erweiterung (GY)

\vspace{0.5em}

\textbf{WICHTIG:} Diese Marker sind NUR für die Lehrkraft. Auf Schülermaterial erscheinen KEINE Niveaumarkierungen!
\end{differenzierung}

\vspace{1em}

\begin{tabularx}{\textwidth}{|l|X|X|X|}
\hline
& \B{} \textbf{Basis} & \Std{} \textbf{Standard} & \E{} \textbf{Erweiterung} \\
\hline
\textbf{Aufgabe 1} & Wochentage zuordnen mit Hilfe & Selbstständig ordnen & + Sätze bilden \\
\hline
\textbf{Aufgabe 2} & Nur volle Stunden & + halbe Stunden & + Viertelstunden \\
\hline
\textbf{Aufgabe 3} & Stundenplan mit Vorlage & Selbstständig ausfüllen & + Beschreiben \\
\hline
\textbf{Aufgabe 4} & Dialog mit Satzanfängen & Dialog mit Stichworten & Freier Dialog \\
\hline
\end{tabularx}

\vspace{1em}

\textbf{Konkrete Hinweise:}
\begin{itemize}
    \item \B{} SuS mit LRS: Wochentage-Poster als Hilfe, Zeitzugabe (+10 Min.)
    \item \B{} DaZ-SuS: Uhrzeiten-Konzept aus Deutsch aktivieren
    \item \E{} Leistungsstarke: Viertelstunden einführen (\es{y cuarto}, \es{menos cuarto})
    \item \E{} Zusatz: Tageszeiten (\es{por la mañana}, \es{por la tarde})
\end{itemize}

% ═══════════════════════════════════════════════════════════════
% 3. HÄUFIGE FEHLER
% ═══════════════════════════════════════════════════════════════

\section*{3. Häufige Fehler}

\begin{fehlerbox}
\begin{tabularx}{\textwidth}{|l|X|X|}
\hline
\textbf{Fehler} & \textbf{Beispiel} & \textbf{Reaktion} \\
\hline
``Halb neun'' falsch & *\es{media nueve} statt \es{ocho y media} & Visualisierung: 8:30 = 8 + 30 Min. \\
\hline
\es{Es la} vs. \es{Son las} & *\es{Son la una} & Regel: 1 = Singular, 2+ = Plural \\
\hline
Miércoles-Akzent & *miercoles & mi-\textbf{ér}-co-les, Akzent betonen \\
\hline
Präposition vergessen & *\es{las ocho} statt \es{a las ocho} & Formel: um = \es{a las} + Zahl \\
\hline
Wochentage Reihenfolge & Durcheinander & Merkspruch, häufige Wiederholung \\
\hline
\end{tabularx}
\end{fehlerbox}

% ═══════════════════════════════════════════════════════════════
% 4. MATERIAL-CHECKLISTE
% ═══════════════════════════════════════════════════════════════

\section*{4. Material-Checkliste}

\begin{materialbox}
\begin{itemize}[label=$\square$]
    \item AB-\sequenz\block.pdf (\lerngruppengroesse{} Kopien)
    \item ML-\sequenz\block.pdf (1× für Lehrkraft)
    \item Lehrwerk: Qué pasa 1, S. 48--53
    \item Kalender (für Einstieg)
    \item Große Uhr mit beweglichen Zeigern
    \item Bildkarten: Wochentage (7 Stück)
    \item Optional: Audio-Track (Wochentage-Lied)
    \item Optional: LP-\sequenz\block.html (Tablets/PCs)
\end{itemize}
\end{materialbox}

% ═══════════════════════════════════════════════════════════════
% 5. LÖSUNGEN
% ═══════════════════════════════════════════════════════════════

\section*{5. Lösungen AB-\sequenz\block}

\textbf{Merkkasten Wochentage:}

\begin{tabular}{lllllll}
L & M & X & J & V & S & D \\
\loesung{Montag} & \loesung{Dienstag} & \loesung{Mittwoch} & \loesung{Donnerstag} & \loesung{Freitag} & \loesung{Samstag} & \loesung{Sonntag} \\
\end{tabular}

\vspace{0.5em}

\textbf{Merkkasten Uhrzeiten:}
\begin{itemize}
    \item \es{Son las ocho y media.} = Es ist halb \loesung{neun}.
\end{itemize}

\vspace{1em}

\textbf{Aufgabe 1: Wochentage ordnen} (4 Punkte)

\begin{tabular}{ccccccc}
1 & 2 & 3 & 4 & 5 & 6 & 7 \\
\loesung{lunes} & \loesung{martes} & \loesung{miércoles} & \loesung{jueves} & \loesung{viernes} & \loesung{sábado} & \loesung{domingo} \\
\end{tabular}

\vspace{1em}

\textbf{Aufgabe 2: Uhrzeiten schreiben} (4 Punkte)

\begin{tabular}{ll}
a) 8:00 & \loesung{Son las ocho.} \\
b) 10:30 & \loesung{Son las diez y media.} \\
c) 1:00 & \loesung{Es la una.} \\
d) 3:30 & \loesung{Son las tres y media.} \\
\end{tabular}

\vspace{1em}

\textbf{Aufgabe 3: Stundenplan} (6 Punkte)

\textit{Individuelle Lösung -- Bewertung:}
\begin{itemize}
    \item Korrekte Fächer auf Spanisch: 2P
    \item Korrekte Wochentage: 2P
    \item Korrekte Uhrzeiten mit \es{a las}: 2P
\end{itemize}

\vspace{1em}

\textbf{Aufgabe 4: Dialog} (6 Punkte)

Beispiellösung:
\begin{itemize}
    \item A: \loesung{¿Cuándo tienes matemáticas?}
    \item B: \loesung{Los lunes a las ocho.}
    \item A: \loesung{¿Y español?}
    \item B: \loesung{Los martes a las diez.}
    \item A: \loesung{¿A qué hora tienes educación física?}
    \item B: \loesung{Los viernes a las once.}
\end{itemize}

\textit{Bewertung: Korrekte Fragestrukturen (3P), passende Antworten mit Tag + Uhrzeit (3P)}

% ═══════════════════════════════════════════════════════════════
% 6. HAUSAUFGABE
% ═══════════════════════════════════════════════════════════════

\section*{6. Hausaufgabe}

\begin{itemize}
    \item Lehrwerk S. 50, Übung 6 (Uhrzeiten schreiben)
    \item Vokabeln lernen: 7 Wochentage + Uhrzeiten
    \item Optional: Lernpfad LP-\sequenz\block.html (interaktive Übungen)
\end{itemize}

% ═══════════════════════════════════════════════════════════════
% 7. REFLEXIONSNOTIZEN
% ═══════════════════════════════════════════════════════════════

\section*{7. Reflexionsnotizen (nach Durchführung)}

\begin{tabularx}{\textwidth}{|l|X|}
\hline
\textbf{Was lief gut?} & \\[2em]
\hline
\textbf{Was lief nicht?} & \\[2em]
\hline
\textbf{Zeitmanagement} & \\[2em]
\hline
\textbf{Für nächstes Mal} & \\[2em]
\hline
\end{tabularx}

\end{document}
